\chapter{Interpretive Content}
\label{ch:philosophy}

Nothing in this chapter is a mathematical result. No claim made below carries \status{ESTABLISHED}, and an occurrence of that tag anywhere in this chapter should be read as an error rather than as a strong statement. What the chapter does instead is say what the formalism of the preceding chapters does and does not license one to believe about it. Declared readings carry \status{DEFINITION}, because a reading is a choice and nothing is proved by making it. Interpretive commitments that the mathematics of earlier chapters actually consumes carry \status{HYPOTHESIS}. There is exactly one such commitment, the observational closure declared in Section~\ref{sec:phil-participatory}. Its tag recurs wherever that same commitment is restated or wherever a later statement inherits it, so the tag occurs more often than the commitment does, and each occurrence names the commitment it carries. Readings whose consequences are unsettled carry \status{OPEN}, with the obstruction named. Where a position is merely available rather than supported, the text says that it is available and says what would be needed to support it, and the difference between availability and support is the distinction this chapter exists to police.

A chapter of this kind is worth having only if the interpretive choices are not idle, and here they are not. One of them is a hypothesis that Chapter~\ref{ch:renormalization} uses in a proposition. Two of them make demands on the geometry that are in direct conflict, in the precise sense that the regime one requires is the regime the other forbids. And one widely used word names two conditions that play opposite roles in Chapters~\ref{ch:coarsegraining} and~\ref{ch:renormalization}, so that conflating them would mistake an operator's degenerate mode for the criterion governing when the operator applies. Sorting this out is work the mathematics does not do for itself, and leaving it undone invites a reader to supply from outside whichever reading they arrived with.

\section{What the base is, and what it is not}
\label{sec:phil-base}

Probability enters this construction through the associated law objects of Chapter~\ref{ch:geometry} and through the finite recognition law of Chapter~\ref{ch:probability}, and nowhere else. An agent is a section-bearing object, so its state channel is
\begin{equation}
q_i\in\Gamma\!\left(\mathcal C_i,
\mathcal E_{\mathrm{state}}\vert_{\mathcal C_i}\right),
\qquad
c\longmapsto q_i(c)\in
\left(\mathcal E_{\mathrm{state}}\right)_c,
\label{eq:phil-fiber-measures}
\end{equation}
whose value at each context is a probability law in the fiber there. Only after choosing a local trivialization is that value represented by a point of the declared model fiber $\mathcal B_{\mathrm{state}}$. The contextual base is the index of that assignment. It carries no probability measure of its own at any point in the development: the design $D=\{c_a\}$ is a declared finite subset rather than a sample from any law, no chapter forms an expectation over contexts, and Open~3.15 records that the construction supplies no probability measure even on a space of sections over $\mathcal C$. Beliefs are laws on the fibers, indexed by base points. \status{DEFINITION} Nothing is proved by saying so; this paragraph restates the typing of Definitions 2.4 and 2.7 and of Chapter~\ref{ch:probability}, and it is recorded here because the readings below turn on it.

One misreading is worth foreclosing before it starts. That agent $i$'s own belief in \eqref{eq:phil-fiber-measures} is defined only on $\mathcal C_i$ should not be confused with the separate fact that a common frame may be unavailable. Proposition 2.22 establishes a global continuous section for the \emph{Gaussian} associated bundle because its selected smooth fiber is contractible. Chapter~\ref{ch:expfamily} now begins more generally, with declared law and kernel fibers that need not be manifolds or contractible; Proposition 2.22 supplies no section theorem for all of them. What Hypothesis 2.10 and its failure control in every case is a principal-bundle trivialization, an identification of fibers with one model fiber, not the existence of a section. In the Gaussian specialization, laws are globally sectionable in this precise sense while a global frame need not be; in a general declared fiber, section existence remains a separate question in the declared category. \status{DEFINITION}

Two readings of that arrangement compete, and neither is preferred by anything in the mathematics.

On the first, the base is a form of intuition in Kant's sense: the a priori structure under which appearances are ordered, rather than the object standing behind them \citep{Kant1781}. What is inferred on this reading is the latent state, the variables $k_i$ and $m_i$ that the generative law of Chapter~\ref{ch:generative} governs and that the recognition law approximates. The base is scaffolding. It exists so that the transports of Chapter~\ref{ch:geometry} are well typed and so that an agent's laws can be indexed at all, and it is not among the things anyone is trying to find out about. \status{DEFINITION}

On the second, the base together with its bundle is noumenal: it is what the representations are representations of, and the latents are the agent-side proxies through which it is approached but never reached. \status{DEFINITION}

The absence of a measure on the base is the first thing a reader will reach for in adjudicating between these, and it does not do the work it appears to do. On the noumenal reading the absence is not an embarrassment but a prediction, and this is the strongest single point in that reading's favor: a noumenon is precisely what one cannot have credences over, so a formalism that equipped the base with a probability would be misdescribing what it had posited. The point must be stated with its limit attached, however, because the competing reading predicts the same absence for its own reason, scaffolding being no more a candidate for credences than a thing-in-itself is. A prediction shared by two hypotheses discriminates neither. The silence of the construction about measures on $\mathcal C$ is therefore evidence-neutral rather than evidence for either party, and the reason the next section is needed is that a criterion which does discriminate has to come from somewhere else. \status{DEFINITION}

Neither reading changes an equation: every proposition of Chapters~\ref{ch:geometry} through~\ref{ch:obstructions} is provable, and every hypothesis statable, under either. That does not make the choice empty, but it does mean the choice cannot be made by checking a derivation, and a reader who wants it settled must accept some criterion the derivations do not supply.

\section{The condition under which a noumenal reading earns its keep}
\label{sec:phil-noumenon}

The criterion this chapter adopts is that a posit which leaves no trace in anything the construction can compute is an idle wheel, and that parsimony removes idle wheels. \status{DEFINITION} This is a declared standard and not a theorem. It is familiar in the form that belief should extend only as far as empirical adequacy requires \citep{vanFraassen1980}, a reader who declines it will reach a different verdict in this section, and nothing in the mathematics forces it on anyone. It is stated openly because the section's conclusion is worthless without it.

Within the graph-link observables actually constructed here, the flat link regime gives the noumenal reading no work to do. Hypothesis 2.10 posits a single reference trivialization over the union of the agents' supports, and Hypothesis 2.14 sets every interaction link equal to the pointwise comparison $\Omega_{ij}=U_iU_j^{-1}$. Proposition 2.11 then returns the identity for the ordered product of comparisons around any closed walk at a common context, and Proposition 2.15 shows that in trivialized coordinates every transported residual becomes a plain difference. This proves that the Regime-I \emph{graph-link sector} carries no nontrivial loop invariant. It does not prove that the base is topologically trivial from graph data---Hypothesis 2.10 already assumed that---and it does not prove that every principal connection on the base is flat. Proposition 2.32 explicitly exhibits a curved, partition-dependent base connection while the local frame-induced pieces remain flat. The idle-wheel verdict is therefore limited to the present population-level graph observables: within that vocabulary the intuitional reading is more economical, but this is the declared criterion's verdict and not a geometric theorem. \status{DEFINITION}

Under nontrivial Regime-II graph-link holonomy the verdict relative to that limited vocabulary reverses, and it reverses for a reason worth stating slowly. In Regime~II the link is declared independently of the frames, and the ordered product around a closed walk,
\begin{equation}
\operatorname{Hol}(\gamma)=\prod_{a=0}^{r-1}L_{i_ai_{a+1}},
\qquad
\operatorname{Hol}(\gamma)\longmapsto g_{i_0}\operatorname{Hol}(\gamma)g_{i_0}^{-1},
\label{eq:phil-holonomy}
\end{equation}
has a gauge-invariant conjugacy class by Proposition 2.13, and by the loop statement of Proposition 2.15 its represented form is conjugate to the ordered product of the trivialized twists $H_{ij}$. Every ingredient belongs to the declared interaction complex. On a connected cluster, the coarsenability theorem of Chapter~\ref{ch:coarsegraining} proves that all such loop products are the identity exactly when the graph-link assignment is a vertex coboundary. Thus a nonidentity product is a model-internal signature of nontrivial \emph{graph-link} structure. It is not, without further data, a signature of principal-bundle topology, a smooth connection, or a curve in the base.

The conclusion under the idle-wheel standard is therefore narrower: the noumenal reading acquires a model-internal referent when the declared graph-link structure has nontrivial holonomy. \status{DEFINITION} The regime is a property of the declared model rather than of the reader, but the conclusion remains an interpretive use of that property. It says neither that the principal bundle is topologically nontrivial nor that the graph invariant arose from base transport. Two qualifications keep those distinct claims out of the conclusion.

The first concerns what the invariant in Equation~\eqref{eq:phil-holonomy} is a signature \emph{of}. Definition 2.12 declares the links on a finite interaction complex that carries no smooth structure and no declared relation to $\mathcal C$, and Open 2.16 records that whether the link holonomy can be realized as the parallel transport of a principal connection along a corresponding closed curve in the base is not settled. Until it is, the invariant detects only the declared graph-link structure, and any inference to base curvature or bundle topology is unlicensed. \status{OPEN} What would close the transport part is exactly what Chapter~\ref{ch:geometry} names: a declared embedding of the interaction complex into $\mathcal C$ sending edges to piecewise-smooth curves, together with a proof that the ordered link product equals the path-ordered exponential of a specified connection along each concatenated loop. Even that equality would establish a connection-holonomy statement, not by itself a theorem that the underlying principal bundle is topologically nontrivial; flat connections can have global holonomy, and nonzero curvature can live on a trivial bundle.

The second is now a split verdict rather than a wholesale exclusion. Chapter~\ref{ch:coarsegraining} proves a theorem-level graph criterion in Regime~II: a connected cluster with positive-definite internal weights is coarsenable in the full fixed-$K$ sense exactly when its internal graph holonomy is trivial, equivalently when its link assignment is a vertex coboundary. What remains open is not that criterion. It is the external choice of a partition and scale, together with an exact gauge-covariant rule for coarse links across cuts and closure of the resulting nonflat family under further merges. Chapter~\ref{ch:renormalization}'s iterated flow is therefore still scoped to its declared flat-link theory space, even though Chapter~\ref{ch:coarsegraining} has already made nontrivial graph holonomy mathematically operative as an admissibility obstruction.

Chapter~\ref{ch:geometry} now distinguishes a third holonomy that this section's argument did not use: the base holonomy of the partition-dependent principal connections constructed in Propositions 2.29 through 2.32. It is logically independent of the graph holonomy of Equation~\eqref{eq:phil-holonomy}. Proposition 2.32 gives an explicit curved average with nonidentity base holonomy and an explicit disagreeing family whose every average is flat; the result depends on smooth variation over the base and on the partition of unity, not on a graph loop. The conjecture worth naming, and not acting on further, is that the noumenal reading could instead be grounded in a specified base connection whose holonomy has an operational population-level trace. \status{CONJECTURE} The curvature example proves mathematical possibility but does not supply that trace, select a canonical connection, or relate it to $L_{ij}$. Settling the conjecture positively requires a declared connection or a partition-independence theorem, an observable constructed from population data, and a proof that the observable changes with the connection's holonomy. Open 2.16 supplies one possible bridge if links are defined as parallel transports along an embedded interaction complex. A proof that every admissible population observable is insensitive to the base connection would instead close the conjecture negatively.

\section{Which realism this is}
\label{sec:phil-realism}

The desktop analogy, on which what an agent perceives stands to what there is as icons stand to machine code, belongs to the interface theory of perception of Hoffman, Singh and Prakash \citep{Hoffman2015,Hoffman2019}. Their argument is evolutionary: perceptual strategies tuned to fitness drive truth-tracking strategies to extinction except where fitness varies monotonically with truth, and the conclusion drawn is that perception is systematically non-veridical, with the representational space related to the world by an arbitrary measurable map carrying no structural constraint whatever.

This construction is doing something different, and the difference is not a matter of emphasis. Agents here are not maps from a world into a private space; by Definition 2.7 they are sections of associated law objects over the base. That typing supports a statement the interface picture cannot make. Hypothesis 2.10 requires the principal bundle to be trivial over the union of the agents' supports, and says exactly what it excludes: the absence of a single smooth reference section to which all local frames can be referred. Read in the other direction, bundle topology constrains the existence of that global reference framing. It does not constrain the existence of each agent's local frame, which local triviality supplies, and a pointwise graph coboundary does not recover the missing smooth \v{C}ech zero-cochain. An arbitrary measurable map admits no analogue of this distinction. Thus the formalism has room for a topological constraint on global representational coordination, while no theorem here identifies graph-link holonomy with that constraint. \status{DEFINITION}

Two restraints belong immediately beside that. No explicit nontrivial probabilistic or renormalization-group realization is developed here: the geometry chapter does analyze non-coboundary \v{C}ech data and its obstruction class, while the tractable probability and flow constructions later impose stronger trivializing hypotheses. And nothing above bears on the veridicality of any individual agent's belief, which is the subject interface theory is actually about; the present construction is compatible with every agent being systematically wrong about everything it represents, and it contains no evolutionary argument of any kind. Nothing here refutes the interface theory, and nothing here is supported by it.

What the mathematics does exhibit is a separation between quantities that survive a change of frame and quantities that do not. On the surviving side stand the conjugacy class of the \emph{graph-link} holonomy of Equation~\eqref{eq:phil-holonomy}, the generalized spectrum of the pair consisting of the Laplacian part and the full interaction operator, and the rank of the coupling matrix,
\begin{equation}
\begin{aligned}
&\bigl[\operatorname{Hol}(\gamma)\bigr]\in\operatorname{Conj}(G),
\qquad
\{d_a\}\ \text{of a declared regular finite pencil }(L,\Lambda),\\
&r=\operatorname{rank}M\in\{0,1,\dots,K\},
\end{aligned}
\label{eq:phil-invariants}
\end{equation}
invariant respectively by Propositions 2.13, 11.9 and 11.11. The spectral entry is conditional: a regular pencil has a finite generalized spectrum, but at the singular fixed operator ray there is no unconditional list of \(NK\) roots. Quotient, mass, and pinning constructions are distinct declared repairs, and an exponent belongs only to a declared thermodynamic or hierarchical family with a limiting counting function. On the other side stand an agent's frame, which a global relabeling moves; the coarse frame, which by Chapter~\ref{ch:coarsegraining} occupies a $K^{2}$-dimensional orbit along which nothing observable varies; and the coupling matrix beyond its rank, since congruence orbits of a positive semidefinite matrix are indexed by rank alone. What the construction licenses one to say, then, is that structure is recoverable from it and intrinsic nature is not.

That is structural realism in the epistemic sense given the position by Worrall, on which what survives is structure rather than nature \citep{Worrall1989}, and it is what this document's reading declares itself to be. \status{DEFINITION} The stronger ontic version, on which structure is not merely what can be known but all that there is, is argued for by Ladyman and collaborators \citep{Ladyman2007,Ladyman2014}, and a moderate form due to Esfeld and Lam places objects and relations on the same footing, individuating objects only relationally and observing that the fiber-bundle formulation is where such a position is most naturally expressed \citep{esfeld2008moderate}. That last observation describes the present setting closely enough to be tempting. It should nonetheless be declined here, because the chapters above exhibit invariants and exhibit unidentifiable quantities, and say nothing whatever about whether the unidentifiable quantities exist. The ontic reading is available and is not supported. \status{OPEN} What would support it is an argument that the unidentifiable directions are absent rather than merely unidentifiable, and no chapter attempts one; an unidentifiability result and an elimination result are different results, and only the first is proved.

\section{Two agreements, and the cost of conflating them}
\label{sec:phil-agreements}

Two conditions on a population are routinely called agreement. They are conditions on different objects, they play opposite roles in the mathematics, and the ambiguity of the single word that names both is expensive. Write them as
\begin{equation}
\begin{aligned}
\text{common graph trivialization:}\quad&
L_{ij}=U_iU_j^{-1}\quad\text{for every edge},\\
\text{belief agreement:}\quad&
z_i=z_j\quad\text{for all }i,j,
\end{aligned}
\label{eq:phil-two-agreements}
\end{equation}
the first a condition on the frames and the links, the second a condition on the trivialized means. \status{DEFINITION}

The first is the coarse-graining admissibility criterion, and its status is no longer conjectural. Chapter~\ref{ch:coarsegraining} proves for a connected cluster with $W_{ij}\succ0$ that trivial internal graph holonomy, existence of vertex elements that trivialize every internal twist, the full $K$-dimensional null space, and annihilation of every internal edge under aggregation are equivalent. Thus the criterion is theorem-level: a cluster satisfying it admits a common graph trivialization even when the population as a whole does not. It is a statement about the declared graph-link assignment, not about a principal connection or the topology of $P$. What remains open is which admissible clusters and scale to choose, how to assign exact gauge-covariant coarse links across cuts when their twists disagree, and whether the enlarged nonflat family is closed under another merge.

The second is a null direction. Under the flat hypothesis the trivialized energy is a sum of self terms and of pair terms in $z_i-z_j$, so the interaction part annihilates the configurations in Equation~\eqref{eq:phil-two-agreements}; Chapter~\ref{ch:coarsegraining}'s translation-invariance theorem states this as one of four equivalent characterizations of the declared family, which is to say that being annihilated on belief agreement is part of what the family \emph{is}. Chapter~\ref{ch:renormalization} then records that at the fixed ray of the declared flow the self sector has vanished, leaving a pure matrix-weighted Laplacian whose kernel contains that configuration; and Chapter~\ref{ch:obstructions} finds the same configuration twice more, as the kernel of a reciprocal pair independently of the link covariances, and as the direction on which a declared apex prior charges exactly $v^{\top}P_0v$. Belief agreement is the degenerate mode. It is what an interaction term costs nothing for, and what a declared self term or prior must charge for if anything is going to.

The two conditions are independent in both directions. A common graph trivialization constrains the vertex elements and the links and leaves the means entirely free, since setting $L_{ij}=U_iU_j^{-1}$ says nothing about $z$ and every configuration of $z$ remains admissible under it. It is not, by itself, the smooth \v{C}ech trivialization of Proposition 2.19. Conversely, transported belief agreement is possible under nontrivial graph holonomy, but only on a subspace: requiring $\mu_i=\rho_k(L_{ij})\mu_j$ along every edge of a closed walk forces
\begin{equation}
\mu_{i_0}\in\ker\bigl(\rho_k(\operatorname{Hol}(\gamma))-I\bigr),
\label{eq:phil-agreement-kernel}
\end{equation}
which is Proposition 12.1 read as a statement about configurations rather than about a precision. Non-trivial holonomy restricts the belief-agreement set to a subspace without emptying it, and trivial holonomy leaves the set whole without entailing that any two agents occupy it. The two conditions constrain one another and neither implies the other.

Because the word covers both, this document avoids ``consensus'' entirely, using ``common graph trivialization'' for the first condition and reserving ``belief agreement'' for the second. \status{DEFINITION} The terminological rule matters because the philosophical temptation runs hard toward the conflation. That coarse agents form when agents agree is an appealing sentence; it is correct on the first reading and wrong on the second, and wrong in a specific and expensive way. Aggregation annihilates internal edges by the algebra of the aggregation matrix regardless of what any agent believes, so belief agreement is not a precondition for forming a coarse agent, and a rule that waited for it would be waiting on a quantity the operation never consults. The deeper error is that belief agreement is the null direction of the very operator being coarse-grained, and that at the fixed ray of the flow the self sector which could charge for it has gone to zero. A criterion built on belief agreement would therefore be reading an operator's degenerate mode as the condition governing when to apply the operator. Common graph trivialization supplies the declared sufficient domain for the parameter recursion and, under positive-definite represented weights, is necessary for full fixed-\(K\) annihilation; semidefinite null-direction exceptions require an effective-support analysis. Belief agreement is instead the direction along which the coarse operator at the flow's fixed ray fails to be definite. \status{DEFINITION} They should not share a name.

\section{The participatory reading and its exact cost}
\label{sec:phil-participatory}

If an agent's observations are the states of other agents, then what an agent directly models is the remainder of the population. \status{HYPOTHESIS} This observational-closure statement is the one interpretive commitment in the chapter that the mathematics uses rather than merely tolerates. Chapter~\ref{ch:renormalization} names it, in the form that an observation between two agents is an edge of the interaction graph, as the hypothesis on which its disconnection argument rests. The graph decomposition is unconditional under that hypothesis. Neither arrival at a fixed ray nor identification of such a ray with a physical law is part of the hypothesis. Arrival additionally requires convergence, uniqueness in the relevant basin, and scheme independence, all recorded as open in Chapter~\ref{ch:renormalization}; the physical-law reading is the separate open empirical and interpretive proposal below. \status{OPEN} This chapter therefore never promotes a fixed point to a physical law by theorem. Observational closure excludes models whose \emph{direct data} have a source outside the population; it does not exclude an unobserved common cause acting on that population.

What the reading buys, beyond the disconnection result, is supplied by Chapter~\ref{ch:obstructions} rather than asserted here. At a variational fixed point of the apex construction, the apex factor's mean is a transported pooling of every constituent's current belief and each constituent's effective prior is that pooling transported back down, so the loop is closed and it lives in the inference, on a joint that remains fixed, acyclic and normalized throughout. The participatory reading is the interpretation under which that loop is the content of the construction rather than a curiosity of its fixed-point equations. Its resemblance to Wheeler's participatory proposal \citep{Wheeler1990} is thematic; nothing here is derived from that proposal and nothing here supports it.

The cost is exact and should be stated without softening, and stating it exactly means separating three commitments that travel together in the usual slogan and do not entail one another. The first is observational closure: every datum any agent has is the state of another agent. That is the commitment named above, it is the only one the disconnection argument of Chapter~\ref{ch:renormalization} consumes, and it is a restriction on the model. \status{HYPOTHESIS} The second is the identification of a repaired limiting object with a physical law. It is logically independent of the first and remains empirically and interpretively inconclusive: no target system, validation criterion, or discriminating evidence is supplied. \status{OPEN} The third is ontological closure: that there is nothing behind the appearances at all. \status{NOT-CLAIMED} The third does not follow from the first two, and the reason is elementary rather than subtle. An unobserved environment may act on every agent alike, entering the effective couplings as a common influence, without ever being the argument of any agent's belief; observational closure constrains what a population directly represents, not what can exist or act on it. This document therefore asserts the first where it is used, leaves the second open, and declines the third. On the first alone, the directly modeled objects are other agents; no conclusion about base topology, base curvature, or external causes follows. The participatory and noumenal readings may pull in different directions under the idle-wheel criterion, but observational closure by itself does not make them logically inconsistent.

The conflict is sharper than a clash of temperaments, because the present versions of the readings use regimes the other does not yet support. Section~\ref{sec:phil-noumenon} gives the noumenal reading a model-internal referent when the Regime-II graph-link structure has nontrivial holonomy. Chapter~\ref{ch:renormalization}'s iterated parameter flow is defined on its declared Regime-I theory space. Chapter~\ref{ch:coarsegraining} proves that trivial internal holonomy is a clean sufficient cluster-level domain and, with positive-definite internal weights, the exact criterion for full fixed-\(K\) annihilation. With semidefinite weights, nontrivial graph holonomy may survive entirely on invisible null directions. The unresolved step therefore includes effective-support typing as well as the cross-cluster link rule, closure under repeated nonflat blocking, and external scale selection. \status{OPEN} One possible dissolution is a closed enlarged family of graph-link operators with an explicit covariant coarse-link or quotient rule. A logically separate dissolution is the conjecture that the noumenal reading rests on an operational trace of a specified base connection, which Proposition 2.32 shows can have nontrivial holonomy even when Regime-I graph loops telescope. The first requires a nonflat graph-link flow; the second requires an observable bridge to base transport. Neither is established, and neither turns a graph link into a principal connection by terminology.

This chapter does not adjudicate. A reader who takes the graph-link noumenal reading gets a model-internal referent for the first invariant of Equation~\eqref{eq:phil-invariants}, but no theorem identifies it with the base; an iterated nonflat or effective-support parameter flow is also not yet defined. A reader who takes the participatory reading gets the component decomposition and the inference-level loop, but no theorem that a physical law is reached. Chapter~\ref{ch:renormalization} records the same refusal for the narrower fixed-point question, leaving the physical-law identification open and empirically inconclusive. The formalism does not decide between the readings, and this document claims neither as established.

\section{What would make any of this testable}
\label{sec:phil-testable}

The construction's own accounting of what it forbids has already been given in Chapter~\ref{ch:renormalization}. Under the commitment of Section~\ref{sec:phil-participatory}, an observation exchanged between two agents is an edge, and the declared blocking map preserves graph components. That gives dynamically independent subsystems; it does not prove that either subsystem converges, that a connected component has one limiting ray, or that a limiting ray is a physical law. Claims about occupied universality classes or no observable variation in a reached law require convergence, basin-level uniqueness, a fixed theory space, and scheme robustness. Those premises remain open. \status{OPEN} The graph decomposition carries its mathematical status in Chapter~\ref{ch:renormalization}; observational closure is the interpretive \status{HYPOTHESIS}; a shared reached physical law is \status{NOT-CLAIMED}.

One model-internal question is decidable now. Given the declared Regime-II links, their ordered products around fundamental cycles can be computed, and Chapter~\ref{ch:coarsegraining}'s theorem determines whether a proposed cluster is graph-trivializing. Such a cluster lies in the declared sufficient domain for full-dimensional aggregation and, when its internal weights are positive definite, this test is also necessary. With semidefinite weights, weighted invisibility must be checked on the represented support and graph holonomy alone is not a complete admissibility test. \status{DEFINITION} This is a consistency and admissibility check on model data, not an empirical detection of bundle topology. Base holonomy is different: Propositions 2.31 and 2.32 make it a functional of smooth frame variation and of a chosen partition of unity. No population observable in the present model computes it, and Open 2.16 supplies the only proposed link-to-curve bridge.

A risky cross-scale test can be specified as an \emph{extension}, but the required empirical structure has not been supplied. Fix in advance: a target system and sampling units; a noise model; an estimator taking training data to gauge-typed estimates of $(U_i,A_i,W_{ij},L_{ij})$ with uncertainty; at least two candidate partitions that pass the established internal-holonomy criterion; a blocking and rescaling map for each; a gauge-alignment rule; a held-out statistic such as conditional log score or a specified covariance moment; an acceptance tolerance; and Gaussian-graphical and spectral baselines. Fit only at the fine resolution. Without refitting at the coarse resolution, push the fitted parameters through each preregistered blocking map and predict the held-out coarse statistic. The added empirical hypothesis is that the predictions agree with the held-out statistic within the declared tolerance and agree across admissible blockings after gauge alignment within propagated uncertainty, while meeting the preregistered baseline margin. Failure of either comparison would falsify that added cross-scale hypothesis. \status{OPEN}

This protocol is deliberately stronger than the present manuscript. No target system, estimator, noise law, tolerance, baseline margin, attraction theorem, or scheme-robust exponent has been declared, so the current empirical status is inconclusive and no existing mathematical proposition would be refuted by a failed data fit. The identification of a repaired fixed object with a physical law remains open, not a verified consequence. The interpretive content is therefore constrained by the mathematics and not confirmed by it: graph-link holonomy, bundle topology, base curvature, and representational gauge redundancy are four different objects, and a viable interpretation must say which one its evidence bears on.

\section{Status register}
\label{sec:phil-register}

{\small
\begin{longtable}{p{0.29\textwidth} p{0.14\textwidth} p{0.45\textwidth}}
\caption{Status register for Chapter~\ref{ch:philosophy}. No entry is \textsc{established}, because no entry is a mathematical result; where a row rests on a proposition of an earlier chapter, that proposition is named and carries its own status there. The chapter's summary claim, that the interpretive content is constrained by the mathematics and not confirmed by it, is likewise not a result.}
\label{tab:phil-register}\\
\toprule
\textbf{Reading or commitment} & \textbf{Status} & \textbf{What is owed} \\
\midrule
\endfirsthead
\multicolumn{3}{@{}l}{\emph{Status register for Chapter~\ref{ch:philosophy}, continued.}}\\
\toprule
\textbf{Reading or commitment} & \textbf{Status} & \textbf{What is owed} \\
\midrule
\endhead
\bottomrule
\endlastfoot
No measure on $\mathcal C$; beliefs are fiber laws, Eq.~\eqref{eq:phil-fiber-measures} & DEFINITION & Restates Defs.~2.4 and 2.7 and the typing of Ch.~\ref{ch:probability}. See Open~3.15. \\
Gaussian law bundle is globally sectionable; general fibers require separate analysis & DEFINITION & Prop.~2.22 uses the contractible Gaussian fiber. What fails without Hyp.~2.10 is a principal-bundle trivialization; no universal section theorem is claimed. \\
Base as form of intuition, and base as noumenon, Sec.~\ref{sec:phil-base} & DEFINITION & Two declared readings. Neither changes an equation, and the missing base measure discriminates neither. \\
Idle-wheel criterion, Sec.~\ref{sec:phil-noumenon} & DEFINITION & A declared standard, in the spirit of \citep{vanFraassen1980}. Declining it changes the verdict below. \\
The noumenal reading has a model-internal referent under nontrivial graph-link holonomy & DEFINITION & Conditional on the idle-wheel criterion. Inputs are Props.~2.11, 2.13 and 2.15 and the coarsenability theorem of Ch.~\ref{ch:coarsegraining}. No bundle-topology inference. \\
Nontrivial graph-link holonomy is a signature of base transport or topology & OPEN & Inherited from Open~2.16 for transport; needs an embedded complex and equality with a specified connection's path-ordered exponential. Topological nontriviality would require an additional argument. \\
The noumenal reading could rest on operational base holonomy instead of graph-link holonomy & CONJECTURE & Props.~2.29--2.32 prove mathematical possibilities but the connection is partition dependent. Needs a specified connection and a population-computable observable sensitive to its holonomy; Open~2.16 is one possible bridge. \\
Bundle topology constrains a global reference framing, Sec.~\ref{sec:phil-realism} & DEFINITION & What Hyp.~2.10 excludes. Local frames still exist; a pointwise graph coboundary is not a smooth \v{C}ech trivialization. \\
Structural realism, not interface theory & DEFINITION & Epistemic sense of \citep{Worrall1989}. Nothing here refutes \citep{Hoffman2015,Hoffman2019} or is supported by it. \\
The ontic form of structural realism & OPEN & Available, not supported. Needs the unidentifiable directions shown absent, not merely unidentifiable \citep{Ladyman2007,esfeld2008moderate}. \\
Graph trivialization and belief agreement are distinct, Eqs.~\eqref{eq:phil-two-agreements} and~\eqref{eq:phil-agreement-kernel} & DEFINITION & Terminology, plus Prop.~12.1 for the second direction of independence. ``Consensus'' is avoided; neither condition is identified with base-bundle triviality. \\
Belief agreement is not a coarsening criterion & DEFINITION & Trivial internal graph holonomy gives the declared sufficient domain and is necessary for full fixed-\(K\) annihilation under positive-definite internal weights. Semidefinite cases require weighted invisibility or effective-support analysis; belief agreement remains only a configuration and null direction. \\
Participatory reading: observational closure, every datum an agent has is the state of another agent & HYPOTHESIS & The one commitment the mathematics consumes, used by Cor.~11.15 with Prop.~11.12. Excludes models whose data have a source outside the population. \\
A repaired limiting object is a physical law & OPEN & R15 remains empirically and interpretively inconclusive. Needs a completed probability repair, a validation criterion, a target system, and discriminating evidence. \\
Ontological closure: there is nothing behind the appearances & NOT-CLAIMED & Declined, not refuted, and it does not follow from the other two: an unobserved environment may act on every agent alike, entering the effective couplings without being the argument of any belief. \\
The present participatory flow and graph-holonomy-based noumenal reading are not yet jointly non-idle & OPEN & Ch.~\ref{ch:coarsegraining} gives a sufficient trivializing domain and the positive-definite exact criterion. Remaining routes require effective-support typing, a closed nonflat coarse-link flow, or an operational base-holonomy observable. \\
Graph components are observationally separate & HYPOTHESIS & Conditional on observational closure; the mathematical content is component preservation in Ch.~\ref{ch:renormalization}. It does not imply one reached limit per component. \\
That a connected population has one shared reached effective law, hence no observable variation in it & NOT-CLAIMED & Asserted nowhere in this chapter without the antecedents of Cor.~11.15, namely Open~11.5, Open~11.13 and Open~11.14, none of which is settled. Prop.~11.12 alone gives one dynamically independent subsystem per component, not one reached limit. \\
Preregistered held-out cross-scale prediction & OPEN & Needs a target system, estimator, noise law, admissible blockings, rescaling, gauge alignment, tolerance, and baselines. Failure would falsify the added empirical hypothesis, not the current mathematical propositions. \\
\end{longtable}
}
